\section{Consistency and crackle (KT)}


%Main mathematical results with proofs (possibly relegated to the Appendices).

We will first prove a number of limit theorems about $d_{f(n)}$, that is the $f(n)$-th death time (when ordered increasingly by size). For consistency we are interested in $d_{f(n)} \to 0$ and $nd_{f(n)}\to \infty$.


We will particularly focus on $d_{\alpha(n-1)}$ for some $\alpha\in [0,1]$. Recall that there are only $n-1$ finite death times so $d_{n-1}$ is the largest finite death value. 

%\subsection[Gaussian d_{n-1} to 0]{Gaussian $d_{n-1} \rightarrow 0$}
%Let $S_n$ denote a random sample of $n$ points from a multivariate Gaussian in $\R^d$.  The probability density function of this Gaussian is $c e^{-\frac{\|x\|^2}{2}}$ for some constant $c$. Note this pdf is spherically symmetric. Denote $f(R)=c e^{-\frac{R^2}{2}}$ for this density in terms of the radial coordinate.



\begin{theorem}\label{thm:d_n-1 Gaussian}
Let $S_n$ denote a random sample of $n$ points from a fixed multivariate Gaussian in $\R^d$. Then for the Rips filtration over $S_n$ we have $d_{n-1} \to 0$ as $n\to \infty$.
\end{theorem}
\begin{proof}
We will first prove the result in the case of a standard Gaussian.

Let $S_n$ denote a random sample of $n$ points from a multivariate Gaussian in $\R^d$. 
 The probability density function of this Gaussian is $c e^{-\frac{\|x\|^2}{2}}$ for some constant $c$. Note this pdf is spherically symmetric. Denote $f(R)=c e^{-\frac{R^2}{2}}$ for this density in terms of the radial coordinate.

Let $X(n,r)$ be the event that $\cup_{x\in S_n} B(x, r)$ is connected. We want to show that there is a function $s(n)$ such that $r(n)\rightarrow 0$ as $n \rightarrow \infty$ and $\mathbb{P}(X(n,s(n)))\rightarrow 1$ as $n\rightarrow \infty$. This implies that $d_*\leq r(n)$ must also limit to $0$.

Set $r(n)$ to satisfy $r(n)^d c \log \log n = 1$. By construction $r(n)\rightarrow 0$ as $n \rightarrow \infty$. We will use $r(n)$ to show that the core is covered by balls around the samples with radius $r(n)$.

Cover $B(0, R)$ with a grid with side lengths $g$. Let $\tilde{C}_n^c$ be the event that at least one of the boxes in this grid is empty. Equation 3.1 in the crackle paper states:
$$
	\mathbb{P}(\tilde{C}(n,R,g)) \leq (2g^{-1})^d R^de^{-ng^d f(R)}.
$$

Set $R_n^c=\sqrt{2(\log n -\log \log n)}$, which will represent the core of the point cloud. Note that
$$
	f(R_n^c)= c e^{-\log n+\log\log n}=c\frac{\log n}{n}.
$$

Set $g=\frac{r(n)}{2\sqrt{d}}$ so $g^d=\frac{1}{c \log \log n 2^d d^{d/2}}$.
\begin{align*}
	\mathbb{P}(\tilde{C}(n,R_n^c,r(n)))
	 & \leq 2^{2d}d^{d/2}c \log \log n (2(\log n -\log \log n))^{d/2}e^{-ng^dc\frac{\log n}{n}} \\
	 & \leq K\frac{(\log n)^{d/2}\log\log n }{\exp(\frac{\log n}{2^d d^{d/2} \log \log n})}
\end{align*}
where $K$ is a constant independent of $n$.

For large $n$, $\log n>(d\log\log n) 2^d d^{d/2} \log \log n$ and thus
$$
	\exp\left(\frac{\log n}{2^d d^{d/2} \log \log n}\right)>\exp(d\log\log n)=(\log n)^d.
$$
Thus for large $n$
$$
	\mathbb{P}(\tilde{C}(n,R_n^c,r(n)))\leq K \frac{(\log n)^{d/2}\log\log n }{(\log n)^d}
$$
which goes to $0$ as $n \to \infty$.

Set $R_n^2:=\sqrt{2(\log n+d\log \log n)}$.

Also directly from the crackle paper we have that the probability of there being a point more that $R_n^2:=\sqrt{2(\log n +d\log\log n)}$ goes to zero as $n\to \infty$. This implies that we can assume all points are within $B(0,R_n^2)$ and that every point in $B(0, R_n^c)$ is within $\frac{1}{ c \log \log n}$ of a sample point.
\begin{align*}
	R_n^2-R_n^c
	 & = \sqrt{2(\log n +d\log\log n)}-\sqrt{2(\log n -\log \log n)}                                \\
	 & =\sqrt{2(\log n)}(\sqrt{1+\frac{d\log\log n}{2\log n}}-\sqrt{1- \frac{\log\log n}{2\log n}}) \\
	 & \sim\sqrt{2(\log n)}((1+\frac{d\log\log n}{4\log n})-(1-\frac{\log\log n}{4\log n})          \\
	 & =\sqrt{2}\sqrt{\log n}\frac{(d+1)\log\log n}{4\log n}                                        \\
	 & =\frac{\sqrt{2}(d+1)}{4}\frac{\log\log n}{\sqrt{\log n}}                                     \\
\end{align*}
Let $s(n)=\max\{r(n), (d+1)\frac{\log\log n}{\sqrt{\log n}}\}$. Then by the above arguments $\P(X(n,s(n)))\rightarrow 1$ as $n\rightarrow \infty$. This implies that $0<d_{n-1}<s(n)$. We also have $s(n)\to 0$ as $n\to \infty$. Together, they squeeze $d_{n-1}$ to satisfy $d_{n-1}\to 0$ as $n\to \infty$.

To then extend the result to general multivariate Gaussian, observe that our sample of points $\hat{S}_n$ is the linear transformation $T$ of a set of points $S_n$ of the standard Gaussian. If $\cup_{x\in S_n} B(x, r)$ is connected, then we can infer that $\cup_{x\in S_n} B(T(x), \|T\|r)$ is also connected. This implies that $d_{n-1}(\hat{S}_n)\leq \|T\|d_{n-1}(S_n)$ and thus $d_{n-1}(\hat{S}_n) \to 0$ as $n\to \infty$.
\end{proof}


\begin{corollary}
    For point clouds sampled from multivariate Gaussian distributions we have $d_* \to 0$ as $n\to \infty$.
\end{corollary}

\begin{proof}
Note that $i_*< n-1$ and so $d_*=d_{i_*}<d_{n-1}$ which we know limits to $0$ almost surely from Theorem \ref{thm:d_n-1 Gaussian}.
\end{proof}


\begin{comment}
\subsection{For more general distributions $d_* \to  0$ if we use percentiles}%{Exponential $d_*\to 0$ but only for percentiles}


Let $S_n$ denote. random sample of $n$ points from an exponential distribution in $\R^d$. The probability density function of this Gaussian is $c e^{-\|x\|}$ for some constant $c$. Note this pdf is spherically symmetric. Denote $f(R)=c e^{-R}$ for this density in terms of the radial coordinate.

Set
$$
R_n:=\log n -3\log \log \log n.
$$
This will the radius of the core in which we will show that all points will be connected at a radius $r_n$ that will decrease to $0$ very slowly. Set $g>0$ such that $g^d=\frac{d+1}{c}(\log \log n)^{-2}$.

Let $X(n,r)$ be the event that $\cup_{x\in S_n\cap B(0, R_n^c)} B(x, r)$ is connected. Set $r_n=2g$
%We want to show that there is a function $s(n)$ such that $r(n)\rightarrow 0$ as $n \rightarrow \infty$ and
%$\P(X(n,s(n)))\rightarrow 1$ as $n\rightarrow \infty$. This implies that $d_*\leq r(n)$ must also limit to $0$.

By construction $r(n)\rightarrow 0$ as $n \rightarrow \infty$.

Cover $B(0, R_n)$ with a grid with side lengths $g$. Let $\tilde{C}_n^c$ be the event that at least one of the boxes in this grid is empty. Equation 3.1 in the crackle paper states:
$$
\mathbb{P}(\tilde{C}(n,R,g)) \leq (2g^{-1})^d R^de^{-ng^d f(R)}.
$$
Set $R_n=\log n - 3\log \log \log n$ which will represent the core of point cloud. Note that
$$
f(R_n)= c e^{-\log n+3\log\log\log n}= c \frac{(\log \log n)^3}{n}.
$$
\begin{align*}
\mathbb{P}(\tilde{C}(n,R_n^c,r(n))) 
\hspace*{-.8cm} &  \\
& \leq\frac{2^d c}{d+1} (\log \log n)^2 (\log n - 3\log \log \log n)^d 
\exp\left(-n \frac{d+1}{c}(\log \log n)^{-2}c \frac{(\log \log n)^3}{n}\right) \\
& \leq \frac{2^d c}{d+1} (\log \log n)^2 (\log n)^d 
\exp\left(-(d+1)(\log \log n)^{-2} (\log \log n)^3\right) \\
& \leq  \frac{2^d c}{d+1} (\log \log n)^2 (\log n)^d 
\exp\left(\log (\log n)^{-d-1}\right) \\
& \leq  \frac{2^d c}{d+1} \frac{(\log \log n)^2 (\log n)^d }{(\log n)^{d+1}}   \\
& \to 0 \text{ as } n \to \infty
\end{align*}

%where $K$ is a constant independent of $n$.
%
%For large $n$, $\log n>(d\log\log n) 2^d d^{d/2} \log \log n$ and thus
%$$\exp\left(\frac{\log n}{2^d d^{d/2} \log \log n}\right)>\exp(d\log\log n)=(\log n)^d.$$
%
%Thus for large $n$
%$$\P(\tilde{C}(n,R_n^c,r(n)))\leq K \frac{(\log n)^{d/2}\log\log n }{(\log n)^d}$$
%which goes to $0$ as $n \to \infty$.

Let $Y(n, R_n)$ denote the points with radius larger than $R_n$, so there are $n-Y$ samples in $B(0, R_n)$. We have that at least $n-Y$ deaths have occurred by radius $r(n)$ as all the samples in $B(0, R_n)$ are connected by then.

We have $\mathbb{E}[Y]=O(\log n)$ which implies
$$\mathbb{P}(Y\geq(1-\alpha) n)\to 0
$$
for all $\alpha \in (0,1)$

This implies that $d_{\alpha n} \to 0$ as $n\to \infty$.
%Also directly from the crackle paper we have that the probability of there being a point more that $R_n^2:=\sqrt{2(\log n +d\log\log n)}$ goes to zero as $n\to \infty$.


---------------- new material Feb 2025 -------------------

Highest density region. 

(Standard definition) -- cite Rob's paper

https://robjhyndman.com/publications/computing-and-graphing-highest-density-regions/


Optimal bandwidth rates: \url{https://otexts.com/weird/03-multivariate.html#bandwidth-matrix-selection}

\url{https://paperpile.com/shared/syX5mqD5uQ12~68AseokG6g}z
p40
conditions on page 31 (for density)
\end{comment}

Theorem \ref{thm:d_n-1 Gaussian} does not hold when we sample our points from distributions with fat tails. We will demonstrate this using the exponential distribution. 


\begin{theorem}
Let $S_n$ denote a random sample of $n$ points from an exponential distribution in $\R^d$. Then for all $k$, $d_{n-k} \to \infty$ as $n\to \infty$. (this is a conjecture)
\end{theorem}

\begin{comment}
\begin{proof}
 The probability density function of this Gaussian is $c e^{-\|x\|}$ for some constant $c$. Note this pdf is spherically symmetric. Denote $f(R)=c e^{-R}$ for this density in terms of the radial coordinate.


Set
$$
R_n:=\log n -3\log \log \log n.
$$
This will be the radius of the core in which we will show that all points will be connected at a radius $r_n$ that will decrease to $0$ very slowly. Set $g>0$ such that $g^d=\frac{d+1}{c}(\log \log n)^{-2}$.

Let $X(n,r)$ be the event that $\cup_{x\in S_n\cap B(0, R_n^c)} B(x, r)$ is connected. Set $r_n=2g$
%We want to show that there is a function $s(n)$ such that $r(n)\rightarrow 0$ as $n \rightarrow \infty$ and
%$\P(X(n,s(n)))\rightarrow 1$ as $n\rightarrow \infty$. This implies that $d_*\leq r(n)$ must also limit to $0$.

By construction $r(n)\rightarrow 0$ as $n \rightarrow \infty$.

Cover $B(0, R_n)$ with a grid with side lengths $g$. Let $\tilde{C}_n^c$ be the event that at least one of the boxes in this grid is empty. Equation 3.1 in the crackle paper states:
$$
\mathbb{P}(\tilde{C}(n,R,g)) \leq (2g^{-1})^d R^de^{-ng^d f(R)}.
$$
Set $R_n=\log n - 3\log \log \log n$ which will represent the core of point cloud. Note that
$$
f(R_n)= c e^{-\log n+3\log\log\log n}= c \frac{(\log \log n)^3}{n}.
$$
\begin{align*}
\mathbb{P}(\tilde{C}(n,R_n^c,r(n))) 
\hspace*{-.8cm} &  \\
& \leq\frac{2^d c}{d+1} (\log \log n)^2 (\log n - 3\log \log \log n)^d 
\exp\left(-n \frac{d+1}{c}(\log \log n)^{-2}c \frac{(\log \log n)^3}{n}\right) \\
& \leq \frac{2^d c}{d+1} (\log \log n)^2 (\log n)^d 
\exp\left(-(d+1)(\log \log n)^{-2} (\log \log n)^3\right) \\
& \leq  \frac{2^d c}{d+1} (\log \log n)^2 (\log n)^d 
\exp\left(\log (\log n)^{-d-1}\right) \\
& \leq  \frac{2^d c}{d+1} \frac{(\log \log n)^2 (\log n)^d }{(\log n)^{d+1}}   \\
& \to 0 \text{ as } n \to \infty
\end{align*}

We want to bound from below the number of samples lying in this core.

$$\int e^{-r}r^{d-1}dr = -e^{-r}\sum_{k=0}^{d-1}[((d-1)!)/((d-1-k)!)]r^{d-1-k} + C$$

%where $K$ is a constant independent of $n$.
%
%For large $n$, $\log n>(d\log\log n) 2^d d^{d/2} \log \log n$ and thus
%$$\exp\left(\frac{\log n}{2^d d^{d/2} \log \log n}\right)>\exp(d\log\log n)=(\log n)^d.$$
%
%Thus for large $n$
%$$\P(\tilde{C}(n,R_n^c,r(n)))\leq K \frac{(\log n)^{d/2}\log\log n }{(\log n)^d}$$
%which goes to $0$ as $n \to \infty$.


Let $Y(n, R_n)$ denote the points with radius larger than $R_n$, so there are $n-Y$ samples in $B(0, R_n)$. We have that at least $n-Y$ deaths have occurred by radius $r(n)$ as all the samples in $B(0, R_n)$ are connected by then.

We have $\mathbb{E}[Y]=O(\log n)$ which implies
$$\mathbb{P}(Y\geq(1-\alpha) n)\to 0
$$
for all $\alpha \in (0,1)$

This implies that $d_{\alpha n} \to 0$ as $n\to \infty$.
%Also directly from the crackle paper we have that the probability of there being a point more that $R_n^2:=\sqrt{2(\log n +d\log\log n)}$ goes to zero as $n\to \infty$.

 \end{proof}

\end{comment}

However, all hope is not lost, because we can prove bounds for percentiles of death times tending to zero.

\begin{definition}
    For $S \subseteq \R^d$ and $\epsilon > 0$, the \textbf{external \( \epsilon \)-covering number} \( N_{\text{ext}}(S,\epsilon) \) is defined as:

\[
N_{\text{ext}}(S, \epsilon) = \min \left\{ k \in \mathbb{N} \mid S \subseteq \bigcup_{i=1}^{k} B(x_i, \epsilon) \text{ for some } x_1, \dots, x_k \in \R^d \right\},
\]
\end{definition}

 
\begin{theorem}
Let $f:\R^d \to \R$ be a Lipschitz probability density function. We want to show that $d_{\alpha (n-1)} \to 0$ where $d_{\alpha (n-1)}$ is the $\alpha$ percentile of finite death times.
\end{theorem}

\begin{proof}
Fix $\alpha\in [0,1)$. Choose $c>0$ such that for $L=\{x|f(x)>c\}$ we have $\P(X\notin L)<(1-\alpha)/2$. 
Choose $\epsilon\in (0,c/6)$ and set $L^\epsilon=\{y+v\mid y\in L, \|v\|\leq \epsilon\}$.


\begin{comment}
Take a maximal $(\varepsilon/2)$-separated subset $S = \{x_1, x_2, \ldots, x_n\} \subseteq L^\epsilon$. 


By maximality of $S$, for any $x \in L$, there exists some $x_i \in S$ such that $|x - x_i| < \varepsilon/2$.

Consider the disjoint balls $\{B(x_i, \varepsilon/4) : i = 1, \ldots, n\}$. These are disjoint because for $i \neq j$:
$$|x_i - x_j| \geq \varepsilon/2 > \varepsilon/4 + \varepsilon/4$$

Since each $x_i \in L$, we have:
$$B(x_i, \varepsilon/4) \subseteq L + B(0, \varepsilon/4) := \{x + y : x \in L, |y| \leq \varepsilon/4\}$$

Therefore:
$$n \cdot \omega_d \left(\frac{\varepsilon}{4}\right)^d = \sum_{i=1}^n |B(x_i, \varepsilon/4)| \leq |L + B(0, \varepsilon/4)|$$

where $\omega_d$ is the volume of the unit ball in $\mathbb{R}^d$.


Therefore, the collection $\{B(x_i, \varepsilon/2) : i = 1, \ldots, n\}$ covers $L$.

This gives us:
$$N_{\text{ext}}(L, \varepsilon/2) \leq n$$

Since $N_{\text{ext}}(L, \varepsilon) \leq N_{\text{ext}}(L, \varepsilon/2)$, we have:
$$N_{\text{ext}}(L, \varepsilon) \leq n$$


By definition, this means:
\begin{enumerate}
\item[(i)] $|x_i - x_j| \geq \varepsilon/2$ for all $i \neq j$ (packing condition)
\item[(ii)] $S$ is maximal: cannot add any more points to $S$ while maintaining (i)
\end{enumerate}


\end{comment}

Packing argument bounds the external covering number, and in turn the number of connected components.


Let $p=\P(X\notin L)$. Note that for all $x\in L$ we have $f(x)>c$. Since $f$ has only finitely many local maxima we have that the number of connected components in $L$ is finite. 
%Denote these connected components by $L_1, L_2, \ldots  L_K$. %Consider also the sets $L_i^\epsilon=\{y\in \R^d| \text{ there exists }x\in L_i \text{ such that } \|x-y\|<\epsilon\}$
%We also know that the volumes of the $L_i$ are all bounded.

Let $Y(L,n)$ denote the random variable for the number of our $n$ samples that are not in $L$. We want $Y(L,n)+K<n(1-\alpha)$ for large $n$. As $K$ is finite, it is sufficient to show $Y(L,n)<5n(1-\alpha)/6$ for large $n$.  

$Y(L,n)$ is a binomial distribution with mean $np$ and variance $np(1-p)$. We can use Chebyshev's inequality to bound $\P(Y(L,n)>5n(1-\alpha)/6)$.
\begin{align*}
    \P(Y(L,n)>5n(1-\alpha)/6)&\leq \P(|Y(L,n)-np|>5n(1-\alpha)/6-n(1-\alpha)/2)\\
    &= \P(|Y(L,n)-np|>n(1-\alpha)/3\\
    &\leq \frac{Var(Y)}{(n(1-\alpha)/3)^2}\\
    &\leq \frac{9np(1-p)}{n^2(1-\alpha)^2}\\
    %&\leq \frac{9(1-\alpha)(1-p)/2}{n(1-\alpha)^2}\\
    &\leq \frac{9}{2n(1-\alpha)}
\end{align*}
This clearly goes to zero as $n$ goes to infinity.


Let $r_n =\frac{k\log(n)}{n}$ (for $k$ constant) which is a function that tends to $0$ as $n$ goes to infinity. Let $L_i^\epsilon=\{y+v\mid y\in L_i, \|v\|\leq \epsilon\}$. For $n$ large, and $\epsilon<c/3$ we can use the Lipshitz condition on $f$ to say that for any $x\in L_i^\epsilon$, the value of $f$ over $B(x,r_n)$ is everywhere at least $c/2>0$.

%For $S \subseteq \R^d$ and $\epsilon > 0$, the \textbf{external \( \epsilon \)-covering number} \( N_{\text{ext}}(S,\epsilon) \) is defined as:

%\[
%N_{\text{ext}}(S, \epsilon) = \min \left\{ k \in \mathbb{N} \mid S \subseteq \bigcup_{i=1}^{k} B(x_i, \epsilon) \text{ for some } x_1, \dots, x_k \in \R^d \right\},
%\]

The external covering number is bounded------



Some argument from the Lipschitz condition that the covering number is bounded by $C_2/\epsilon^d$ for some constant $C_2$, for $\epsilon$ sufficiently small.

--------------------------------------------------------------------

Let $Z_i(L,r_n,n)$ be the event that $\cap_{x_i\in L_i^\epsilon} B(x_i, r_n) \cap L_i$ is connected. We want to prove that as $n$ increases we have $\P(Z(L,r_n, n)=1) \to 1$.  


\end{proof}

\begin{comment}
We also want to bound the rate of $d_{\alpha(n-1)}$ from below. To do this we will prove that $n^{1/d}d_{\alpha(n-1)} \to \infty$ as $n\to \infty$.

We will provide this bound by lower bounding the number of isolated vertices (as each of these will contribute their own connected component).


We want to lower bound the expected number of connected components using the fact that the probability density function \( f(x) \) is bounded above by \( K \), i.e., \( f(x) < K \) for all \( x \in \mathbb{R}^d \).  




Let \( f(x) \) be a probability density function over \( \mathbb{R}^d \) that is continuous and satisfies \( f(x) \leq K \). Suppose we generate a set \( S \) of \( n \) i.i.d. points according to \( f(x) \), and we construct a random geometric graph by connecting points within a radius \( r \). 

Let \( X \) be a randomly chosen point from the set \( S \) of \( n \) i.i.d. points drawn from \( f(x) \). The point \( X \) is isolated if no other \( n-1 \) points fall within a ball of radius \( r \) centered at \( X \), denoted as \( B(X, r)\).  

The probability that a single other point lands inside \( B(X, r) \) is  

\[
P(\text{random point falls in } B(X, r) \mid X = x) = \int_{B(x, r)} f(y) dy\leq  K V_d r^d,
\]

where \( V_d \) is the volume of the unit ball in \( \mathbb{R}^d \), given by  

\[
V_d = \frac{\pi^{d/2}}{\Gamma(1 + d/2)}.
\]

Since the remaining \( n-1 \) points are independently drawn from \( f(x) \), the probability that none of them land in \( B(X, r) \) is bounded from below by 
\(
\Bigl(1 -  K V_d r^d \Bigr)^{n-1} \approx \exp\Bigl(- (n-1) K V_d r^d \Bigr),
\)
where we used the approximation \( (1 - u)^{n-1} \approx e^{-(n-1)u} \) for small \( u \).


The expected number on the expected number of isolated nodes is thus bounded below by
\[
\mathbb{E}[\# \text{ isolated vertices}] \geq n \exp\Bigl(-n V_d r^d K\Bigr).
\]

If we set $r$ by $r^d=\frac{-\log(2(1-alpha))}{nV_dK}$ then we have $n \exp\Bigl(-n V_d r^d K\Bigr)=n\exp(\log(1-\alpha))=2(1-\alpha)n$. 



\end{comment}
%\subsection[ndalpha(n-1) to infinity]{$nd_{\alpha(n-1)}\to \infty$}
For consistency we also need to prove lower bounds.


A consequence of the below theorems is that $nd_{\alpha(n-1)}\to \infty$ %whenever $d_*$ is chosen as the any fixed percentile of the death times. 
This holds for very general distributions. We only require a global upper bound on the pdf and the dimension of the space to be at least $2$.

\begin{theorem}
	Let $d_k$ denote the $k$th smallest death time of the $0$-homology for the Rips filtrations of a set of $n$ points sampled from probability distribution $X$ over $\R^2$ such that the density is bounded above. Note that there are $n-1$ finite death times.

	Choose any sequence of integers $\{\omega(n)\}$ such that $0<\omega(n) <n$ for all $n$ and $\lim_{n\to \infty} \omega_n =\infty$. Then $nd_{\omega(n)}\to \infty$ as $n\to \infty$ with probability $1$.
\end{theorem}
\begin{proof}
	It is sufficient for there to exist a function $r(n)$ such that $\mathbb{P}(d_{\omega(n)})\geq \frac{r(n)}{n}) \to 1$ as $n\to \infty$ and $r(n)\to \infty$ as $n\to \infty$.
	Let $C(n, r)$ denote the number of components of $\cup_x B(x, r)$.

	Now $d_{\omega(n)} > \frac{r(n)}{n}$ occurs when $C(n, \frac{r(n)}{n})> n - \omega(n)$.

	Since the number of components can only be an integer between $1$ and $n$,
	$$
		\mathbb{P}\left(C\left(n, \frac{r(n)}{n}\right)> n-\omega(n)\right) \geq \frac{\mathbb{E}[C(n, \frac{r(n)}{n})]-(n-\omega(n))}{n-(n-\omega(n))}= \frac{\mathbb{E}[C(n, \frac{r(n)}{n})]-n+\omega(n)}{\omega(n)}.
	$$
	This implies that it is sufficient to show that there exists a suitable $r(n)$ such that $r(n)\to \infty$ as $n\to \infty$. and
	$$
		\lim_{n\to \infty}\frac{\mathbb{E}[C(n, \frac{r(n)}{n})]-n+\omega(n))}{\omega(n)} = 1.
	$$
	For small radii $r$ a useful lower bound on the number of connected components is $n-E(n,r)$ where $E(n,r)$ is the number of edges in the graph connecting any two points are less that $2r$ apart, that is
	$\mathbb{E}[C(n, \frac{r(n)}{n})]\geq n-\mathbb{E}[E(n,\frac{r(n)}{n})]$. The expected number of edges is much easier to calculate than the expected number of components.
	\begin{align*}
		\mathbb{E}[E(n,\frac{r(n)}{n})]
		 & =\frac{n(n-1)}{2}\mathbb{P}(d(x,y)<2r\mid x,y \sim X) \\
		 & \leq \frac{n(n-1)}{2} K\pi (\frac{r(n)}{n})^2         \\
		 & \leq \frac{K \pi}{2}r(n)^2
	\end{align*}
	where $K$ is the global upper bound on the pdf of $X$.

	Set $r(n)=\log (\omega (n)+1)$. We automatically get $\lim_{n\to \infty} r(n)= \infty$ as $\lim_{n\to \infty} \omega(n)=\infty$.

	$\mathbb{E}[C(n, \frac{r(n)}{n})]\geq n- K \pi\log(\omega(n)+1)$ and thus
	$$
		\lim_{n\to \infty} \frac{\mathbb{E}[C(n, \frac{r(n)}{n})]-n+\omega(n)}{\omega(n))} \geq \lim_{n\to \infty} \frac{\omega(n)-K \pi\log(\omega(n)+1)}{\omega(n)}= 1.
	$$
\end{proof}

\begin{theorem}
	Let $d_1$ denote the smallest death time of the $0$-homology for the Rips filtrations of a set of $n$ points sampled from probability distribution over $\R^d$, $d\geq 3$ such that the density is bounded above.

	Then $nd_{1}\to \infty$ as $n\to \infty$ with probability $1$.
\end{theorem}
\begin{proof}
	The argument here is an easier version of the above theorem. We wish to find a function $r(n)$ such that $\lim_{n\to \infty}r(n)=\infty$ and
	$$
		\lim_{n\to \infty}\mathbb{P}(C(n, \frac{r(n)}{n})=n)=1.
	$$
	Set $r(n)=\log(n)$.
	\begin{align*}
		\mathbb{E}[E(n,\frac{r(n)}{n})]
		 & =\frac{n(n-1)}{2}\mathbb{P}(d(x,y)<2r(n)\mid x,y \sim X) \\
		 & \leq \frac{n(n-1)}{2} K\pi s_d(\frac{2\log(n)}{n})^d     \\
		 & \leq 2^{d-1}Ks_d \frac{(\log(n))^d}{n^{d-1}}
	\end{align*}
	where $K$ is the global upper bound on the pdf of $X$ and $s_d$ is the volume of the unit ball in $\R^d$.

	$\mathbb{E}[C(n, \frac{r(n)}{n})]\geq n- 2^{d-1}Ks_d\frac{(\log(n))^d}{n^{d-1}}$.

	Since the number of connected components must be an integer between $1$ and $n$ we deduce
	\begin{align*}
		\mathbb{P}(C(n, \frac{r(n)}{n})=n)
		 & \geq \mathbb{E}[C(n, \frac{r(n)}{n})]+1-n \geq n-2^{d-1}Ks_d \frac{(\log(n))^d}{n^{d-1}}+1-n \\
		 & =1-2^{d-1}Ks_d \frac{(\log(n))^d}{n^{d-1}}                                                   \\
	\end{align*}
	and thus $\lim_{n\to \infty}\mathbb{P}(C(n, \frac{r(n)}{n})=n)=1.$
\end{proof}

\subsection{Bounds for kernel density estimates}
\begin{lemma}\label{lemma:boundsforkde}
	Let $\{\bm{x}_1, \dots, \bm{x}_n\}$ be $n$ data points in $\mathbb{R}^p$, with kernel density estimate given by
	$$ \hat{f}(\bm{x}) = \frac{1}{n} \sum_{i = 1}^n |\bm{H}|^{-1/2}K\left(\bm{H}^{-1/2}(\bm{x} - \bm{x}_i)\right),
	$$
	where $K(\bm{u})$ is a multivariate kernel function that is bounded below by 0 and above by $K(\bm{0}) = M$, and $\bm{H}$ is a $p\times p$ symmetric positive-definite bandwidth matrix. Then for all $j \in \{1, \dots, n\}$, the kernel density estimates at the observations, $\hat{f}(\bm{x}_j)$, and their negative logarithms,  $- \log \hat{f}(\bm{x}_j)$, are bounded by constants that depend only on $p$ and $n$.
\end{lemma}
\begin{proof}
	$$
		\frac{M|\bm{H}|^{-1/2}}{n} \leq\hat{f}(\bm{x}_j)  = \frac{1}{n} \sum_{i = 1}^n |\bm{H}|^{-1/2}K\left(\bm{H}^{-1/2}(\bm{x}_j - \bm{x}_i)\right) \leq M|\bm{H}|^{-1/2}\, .
	$$
  The constant $M$ may depend on $p$.
\end{proof}

For example, if we use the multivariate Gaussian kernel $K(\bm{u}) = \frac{1}{(2 \pi)^{p/2}} \exp \left( -\frac12 \|\bm{u}\| \right)$, then $M = \frac{1}{(2 \pi)^{p/2}}$.  If we use the multivariate Epanechnikov kernel $K(\bm{u}) = \frac{p+2}{2c_p} \left( 1 - \|\bm{u}\| \right)_{+}$, where $c_p$ is the volume of the unit sphere in $\mathbb{R}^p$, and $u_+ = max(0, u)$, then $M = (p+2)/(2c_p)$.

% When there is a bandwidth $h$, the lower bound of $\hat{f}(\bm{x}_j)$ has an additional $h$ in the denominator.


